\documentclass{beamer}
\usepackage{amsfonts,amsmath,oldgerm,mathrsfs,booktabs}
\usepackage{tikz}
\usetheme{sintef}
%\usepackage{xeCJK}
\linespread{1.5}

\newcommand{\testcolor}[1]{\colorbox{#1}{\textcolor{#1}{test}}~\texttt{#1}}

\def\b{\color{blue}}
\def\r{\color{red}}
\def\d{{\rm d}}

\usefonttheme[onlymath]{serif}

\titlebackground*{assets/background}

\newcommand{\hrefcol}[2]{\textcolor{cyan}{\href{#1}{#2}}}

\title{{\Large The Parametric Complexity of Operator Learning}}

\vspace{5cm}
%\subtitle{subtitle}
% \course{Master's Degree in Computer Science}
\author{Samuel Lanthaler, Andrew M. Stuart}
%\vspace{10cm}
% \supervisor{}
% \IDnumber{1234567}
\date{ \vspace{15pt}
{\large Presenter: Qiqi Rao}\\
\medskip
May 18, 2026
}

\begin{document}
\maketitle


\section{{\Large Motivation and Background}}

\begin{frame}[fragile]{\thesection. Operator Learning}
\framesubtitle{$\quad$}
{\Large {\b What is operator learning?}}
\medskip

{\Large Neural operator architectures employ neural networks to approximate operators $ S^\dagger : X \to Y $ mapping between Banach spaces of functions.}
\medskip

{\Large {\b Applications:}
\begin{itemize}
    \item Accelerate model evaluations via emulation (e.g., solving PDEs)
    \item Discover input--output maps from data when no model is available
\end{itemize}
}
\medskip

{\Large {\b Key architectures:} PCA-Net, DeepONet, Fourier Neural Operator (FNO)}
\medskip

{\Large {\b Central question:} {\r How many parameters are needed to achieve a given approximation accuracy $ \epsilon $?}}

\end{frame}

\begin{frame}{\thesection. The Encoding--Approximation--Reconstruction Framework}
\framesubtitle{$\quad$}
{\Large Operator learning frameworks decompose as:}
\medskip

{\Large 
\begin{enumerate}
    \item {\b Encoding $ E $}: Map infinite-dimensional input $ u \in X $ to finite-dimensional representation $ E(u) \in \mathbb{R}^{D_X} $
    \item {\b Neural Network $ \Psi $}: Map encoded input to encoded output $ \Psi : \mathbb{R}^{D_X} \to \mathbb{R}^{D_Y} $
    \item {\b Reconstruction $ R $}: Reconstruct output function from finite-dimensional representation
\end{enumerate}
}
\medskip

{\Large {\b Neural operator}: $ S = R \circ \Psi \circ E \approx S^\dagger $}
\medskip

{\Large {\b Complexity} $ := $ size of $ \Psi $ (number of nonzero weights and biases)}

\end{frame}

\begin{frame}{\thesection. Curse of Dimensionality for Neural Networks}
\framesubtitle{$\quad$}
{\Large {\b Finite-dimensional CoD} (Yarotsky, 2017):}
\medskip

{\Large For functions $ f \in C^r([0,1]^D) $, any ReLU network $ \Psi $ achieving accuracy $ \epsilon $ must have:
$$
\text{size}(\Psi) \ge \epsilon^{-\gamma D / r}
$$
}
\medskip

{\Large {\b Intuition}: As dimension $ D \to \infty $, the complexity grows exponentially.}
\medskip

{\Large {\b Key insight for operator learning}: Since neural operators map between {\r infinite-dimensional} spaces, the encoded dimension $ D_X, D_Y \gg 1 $, suggesting a ``scaling limit'' of the CoD.}

\end{frame}

\section{{\Large The Curse of Parametric Complexity}}

\begin{frame}{\thesection. Main Objective}
\framesubtitle{$\quad$}
{\Large {\b Two main contributions of this paper:}}
\medskip

{\Large {\b 1.} {\r Prove a general curse}: For operators characterized only by $ C^r $- or Lipschitz-regularity, operator learning suffers from a ``curse of parametric complexity'' --- complexity grows {\r exponentially} in $ \epsilon^{-1} $.}
\medskip

{\Large {\b 2.} {\b Overcome the curse}: For Hamilton--Jacobi (HJ) equations, leverage additional structure (method of characteristics) to achieve {\b algebraic} complexity bounds.}
\medskip

{\Large {\b Key message}: Regularity alone is not enough. One must exploit {\r problem-specific structure} to achieve efficient operator learning.}

\end{frame}

\begin{frame}{\thesection. The Curse of Parametric Complexity (Theorem 2.11)}
\framesubtitle{$\quad$}
{\Large {\b Setting}: Let $ K \subset X $ be compact, containing an infinite-dimensional hypercube $ Q_\alpha $ with decay rate $ \alpha > 1 $.}
\medskip

{\Large {\b Result}: For any $ r \in \mathbb{N} $ and $ \delta > 0 $, there exists a $ C^r $-functional $ S^\dagger : K \to \mathbb{R} $ such that approximation to accuracy $ \epsilon $ requires:
$$
\text{cmplx}(S_\epsilon) \ge \exp\!\big(c \, \epsilon^{-1/(\alpha + 1 + \delta)r}\big)
$$
}
\medskip

{\Large {\b Interpretation}: 
\begin{itemize}
    \item No algebraic rate is achievable for general $ C^r $-regular operators
    \item Applies to PCA-Net, DeepONet, FNO, and NOMAD
    \item This is an infinite-dimensional analogue of the classical CoD
\end{itemize}
}

\end{frame}

\begin{frame}{\thesection. Intuition Behind the Curse}
\framesubtitle{$\quad$}
{\Large {\b Why does the curse arise?}}
\medskip

{\Large 
\begin{itemize}
    \item Compact sets $ K \subset X $ in infinite-dimensional spaces contain ``infinite-dimensional hypercubes'' $ Q_\alpha = \big\{ A \sum_{j=1}^\infty j^{-\alpha} y_j e_j : y_j \in [0,1] \big\} $
    \item As we increase the encoding dimension $ D_X $, we access more and more ``directions'' in $ K $
    \item The finite-dimensional CoD kicks in for each new direction
    \item The scaling limit $ D \to \infty $ leads to {\r exponential} lower bounds
\end{itemize}
}
\medskip

{\Large {\b Consequence}: To beat the curse, one must leverage structure {\r beyond} smoothness.}

\end{frame}

\begin{frame}{\thesection. Sketch of Proof for Theorem 2.11}
\framesubtitle{$\quad$}
{\Large {\b Main idea}: Reduce the infinite-dimensional operator problem to a hard high-dimensional neural network approximation problem.}
\medskip

{\Large
\begin{enumerate}
    \item {\b Embed a cube}: Show that the compact set $ K $ contains an infinite-dimensional cube $ Q_\alpha $, so inputs can be parameterized by coefficients $ (y_1,y_2,\dots) $
    \item {\b Truncate}: Restrict attention to the first $ D $ coordinates; this creates a $ D $-dimensional approximation problem inside the operator problem
    \item {\b Apply finite-dimensional CoD}: Use Yarotsky-type lower bounds for ReLU networks on $ [0,1]^D $
    \item {\b Optimize $ D $ vs.\ $ \epsilon $}: Choosing $ D $ as a function of the target accuracy turns the finite-dimensional lower bound into an {\r exponential} lower bound in $ \epsilon^{-1} $
\end{enumerate}
}
\medskip

{\Large {\b Intuition}: Even though the operator acts on functions, hidden inside the function space there are arbitrarily many hard finite-dimensional directions.}

\end{frame}

\section{{\Large Overcoming the Curse: HJ-Net}}

\begin{frame}{\thesection. The Hamilton--Jacobi Equation}
\framesubtitle{$\quad$}
{\Large {\b HJ equation}: Given initial data $ u_0 $, find $ u(q,t) $:
$$
\partial_t u + H(q, \nabla_q u) = 0, \quad u(t=0) = u_0
$$
on periodic domain $ \Omega = [0, 2\pi]^d $.
}
\medskip

{\Large {\b Solution operator}: $ S_t^\dagger : u_0 \mapsto u(\cdot, t) $}
\medskip

{\Large {\b Method of characteristics}: The solution can be computed via the Hamiltonian ODE system:
$$
\dot{q} = \nabla_p H(q,p), \quad \dot{p} = -\nabla_q H(q,p), \quad p_0 = \nabla_q u_0(q_0)
$$
}
\medskip

{\Large {\b Key observation}: The same Hamiltonian flow map $ \Psi_t^\dagger $ is used for all inputs; once $ (q_0,p_0,z_0) $ is encoded from $ u_0 $, the evolution is determined by the Hamiltonian $ H $.}

\end{frame}

\begin{frame}{\thesection. HJ-Net Architecture}
\framesubtitle{$\quad$}
{\Large {\b Three-step operator}: $ S_t = R \circ P \circ \Psi_t^N \circ E $}
\medskip

{\Large {\b Step (a) Encode}: Evaluate $ u_0 $ at $ N $ points:
$$
E(u_0) = \{(q_0^j, \nabla u_0(q_0^j), u_0(q_0^j))\}_{j=1}^N
$$
}

{\Large {\b Step (b) Flow}: Apply neural network $ \Psi_t \approx \Psi_t^\dagger $ to each encoded triple:
$$
(q_0^j, p_0^j, z_0^j) \mapsto (q_t^j, p_t^j, z_t^j)
$$
}

{\Large {\b Step (c) Reconstruct}: Interpolate $ \{(q_t^j, z_t^j)\}_{j=1}^N $ via moving least squares}
\medskip

{\Large {\b Why this works}: The flow map $ \Psi_t^\dagger : \mathbb{R}^{2d+1} \to \mathbb{R}^{2d+1} $ is a {\r fixed-dimensional} smooth function --- no infinite-dimensional curse!}

\end{frame}

\begin{frame}{\thesection. HJ-Net Beats the Curse (Theorem 5.1)}
\framesubtitle{$\quad$}
{\Large {\b Result}: For HJ equation with $ C^r $ initial data ($ r \ge 2 $) and $ C^{r+1} $ Hamiltonian, the HJ-Net achieves:
$$
\sup_{u_0 \in F} \| S_t(u_0) - S_t^\dagger(u_0) \|_{L^\infty} \le \epsilon
$$
with {\b algebraic} complexity bounds:}
\medskip

{\Large 
\begin{itemize}
    \item Number of encoding points: $ N \le C \epsilon^{-d/r} $
    \item Neural network size: $ \text{size}(\Psi_t) \le C \epsilon^{-(2d+1)/r} \log(\epsilon^{-1}) $
    \item Neural network depth: $ \text{depth}(\Psi_t) \le C \log(\epsilon^{-1}) $
\end{itemize}
}
\medskip

{\Large {\b Key point}: Complexity grows only {\b polynomially} in $ \epsilon^{-1} $ --- the curse is overcome!}

\end{frame}

\begin{frame}{\thesection. Why HJ-Net Succeeds}
\framesubtitle{$\quad$}
{\Large {\b Mechanisms for overcoming the curse:}}
\medskip

{\Large 
\begin{enumerate}
    \item {\b Separation of concerns}: The Hamiltonian flow $ \Psi_t^\dagger $ does not depend on the input $ u_0 $; it is a {\r finite-dimensional} map $ \mathbb{R}^{2d+1} \to \mathbb{R}^{2d+1} $
    \item {\b Structure exploitation}: The method of characteristics provides an exact decomposition of the solution operator
    \item {\b Scattered data interpolation}: Moving least squares provides a stable reconstruction from the transported data points
\end{enumerate}
}
\medskip

{\Large {\b Contrast with generic approaches}: PCA-Net, DeepONet, FNO treat the operator as a ``black box'' and cannot avoid the curse for general operators.}

\end{frame}

\begin{frame}{\thesection. Sketch of Proof for Theorem 5.1}
\framesubtitle{$\quad$}
{\Large {\b Main idea}: Decompose the error according to the three steps of HJ-Net.}
\medskip

{\Large
\begin{enumerate}
    \item {\b Approximate the flow}: The characteristic flow $ \Psi_t^\dagger : \mathbb{R}^{2d+1} \to \mathbb{R}^{2d+1} $ is finite-dimensional and smooth, so a ReLU network approximates it with algebraic cost
    \item {\b Choose enough sample points}: Take $ N $ encoding points so the transported point cloud is dense enough in space
    \item {\b Reconstruct stably}: Moving least squares turns the transported values $ (q_t^j,z_t^j) $ into an approximation of $ u(\cdot,t) $
    \item {\b Combine the bounds}: Add the flow approximation error and reconstruction error, then balance parameters to make both of order $ \epsilon $
\end{enumerate}
}
\medskip

{\Large {\b Key difference from the negative result}: The proof uses special structure of the HJ equation --- characteristics reduce the hard operator problem to learning a manageable finite-dimensional flow.}

\end{frame}

\section{{\Large Conclusion and Future Work}}

\begin{frame}{\thesection. Summary of Contributions}
\framesubtitle{$\quad$}
{\Large {\b Contribution 1 --- The Curse (Negative Result):}
\begin{itemize}
    \item General $ C^r $-regular operators require exponential complexity
    \item Applies to PCA-Net, DeepONet, FNO, NOMAD
    \item Regularity alone is insufficient for efficient operator learning
\end{itemize}
}
\medskip

{\Large {\b Contribution 2 --- Overcoming the Curse (Positive Result):}
\begin{itemize}
    \item HJ-Net exploits the characteristic structure of HJ equations
    \item Achieves algebraic (polynomial) complexity bounds
    \item The flow map is finite-dimensional $\Rightarrow$ no infinite-dimensional curse
\end{itemize}
}
\medskip

{\Large {\b Overall message}: Efficient operator learning requires leveraging {\r problem-specific structure} beyond smoothness.}

\end{frame}

\begin{frame}{\thesection. Limitations}
\framesubtitle{$\quad$}
{\Large 
\begin{enumerate}
    \item {\b Classical solutions only}: The approach based on characteristics requires that classical solutions exist; it cannot handle viscosity solutions after singularity formation
    \item {\b Short time intervals}: The analysis is restricted to $ t \in [0, T] $ where $ T < T^*(F) $, the maximal existence time of classical solutions
    \item {\b Fixed spatial dimension}: The paper does not address the CoD with respect to the spatial dimension $ d $ (only the infinite-dimensional input space)
    \item {\b ReLU networks only}: The lower bounds are established for feedforward ReLU networks; extensions to other activations remain open
    \item {\b No semigroup property}: The approximate operator $ S_t $ does not satisfy the semigroup property, so iteration over time is not directly possible
\end{enumerate}
}

\end{frame}

\begin{frame}{\thesection. Future Work}
\framesubtitle{$\quad$}
{\Large 
\begin{enumerate}
    \item {\b Viscosity solutions}: Extend the methodology beyond classical solutions to handle singularities and shocks in HJ equations
    \item {\b Symplectic neural networks}: Use symplectic architectures to learn the Hamiltonian flow, preserving geometric structure
    \item {\b Combining with CoD w.r.t. $ d $}: Integrate with existing work on beating the curse of dimensionality in the spatial dimension
    \item {\b $ C^1 $ topology estimates}: Generalize error estimates to the $ C^1 $ topology to enable operator iteration over longer time intervals
    \item {\b Information-theoretic perspective}: Characterize the curse independently of activation function (cf. Lanthaler, 2024)
    \item {\b Other PDE classes}: Identify structure in other PDEs that can be leveraged similarly to characteristics in HJ equations
\end{enumerate}
}

\end{frame}

\backmatter
\end{document}
